\section{The same raw-mean obstruction on the unit circle}
\label{sec:raw-student-circle}
The preceding example does not rely on different input norms. A fixed two-dimensional sphere already permits the same divergence.

\begin{corollary}[A unit-norm input counterexample]
\label{cor:raw-student-circle}
The raw-mean obstruction holds for the one-neuron network
\[
 f_{v,w_1,w_2,\gamma}(x)=v\tanh(w_1x_1+w_2x_2+\gamma),
\]
with four independent coordinates of density $q_p$, when every input lies on the unit circle. The prior, architecture and Gaussian likelihood variance remain fixed. There are noiseless bounded-label laws for which the raw posterior mean has squared test risk at least $c n^{1/3}$ with fixed positive probability, while the zero comparator has risk at most $n^{-1/3}$.
\end{corollary}
\begin{proof}
Retain $b=n^{-1/6}$, $R=e^{c n^{2/3}}$, and $a=b/R$ from the scalar construction, with its sufficiently small fixed $c$. Put $r=\sqrt{1-a^2}$. Replace the common inputs by
\[
 x_+=(a,r),\qquad x_-=(-a,r),
\]
with labels $b,-b$ and probabilities $(1-1/n)/2$ each. The rare input is $x_*=(1,0)$, with label zero and probability $1/n$. Every input has norm one. For large $n$, $r\in[1/2,1]$.

On the no-rare, balanced-count event, the training likelihood depends on $w_2,\gamma$ only through $h=rw_2+\gamma$. The joint density of $(h,\gamma)$ is
\[
 j_r(h,\gamma)=q_p(\gamma)q_p((h-\gamma)/r)/r.
\]
Its marginal $q_r(h)=\int j_r(h,\gamma)d\gamma$ is a probability density, is bounded above uniformly in $r\in[1/2,1]$, and is bounded below by a positive constant on $|h|\le1$, uniformly in that range. The upper bound follows from convolution with a bounded density. For the lower bound, integrate over $|\gamma|\le1$ and use the positive minima of $q_p$ on fixed compact intervals.

Consequently the evidence upper bound from the scalar construction applies unchanged, with $h$ replacing the hidden bias: the saturation estimate uses only that its density integrates to one, and the linear-region estimate uses only a uniform density upper bound. Thus
\[
 Z_{S_n}\le C Z_*,\qquad Z_*=R^{-p}/(bn).
\]
The same global sign symmetry permits restriction to $v>0$.

It remains to check the signed numerator at $x_*$. Define
\[
 I_r(w,h)=\int\tanh(w+\gamma)j_r(h,\gamma)d\gamma.
\]
On the good-fit rectangle, $w=w_1>0$ and $|h|<aw$, so $w+h>0$. Pair the integral around $\gamma=-w$. For $t>0$, monotonicity of $q_p$ in absolute value gives
\begin{align*}
 q_p(-w+t)&\ge q_p(-w-t),\\
 q_p((h+w-t)/r)&\ge q_p((h+w+t)/r).
\end{align*}
Both inequalities use positive centers, $w>0$ and $h+w>0$. Therefore
\[
 I_r(w,h)=\int_0^\infty\tanh(t)
       [j_r(h,-w+t)-j_r(h,-w-t)]dt\ge0.
\]
Thus the full good-fit region has no negative contribution after integrating the latent bias decomposition.

Moreover, there are fixed $W_0,c_0>0$ such that
\[
 I_r(w,h)\ge c_0\qquad(w\ge W_0,\ |h|\le1,\ 1/2\le r\le1).
\]
To verify this, use only the paired terms with $\gamma=-w+t\in[-1,1]$. Their positive density $j_r(h,\gamma)$ is bounded below by an absolute positive constant. The reflected density at $-2w-\gamma$ is at most $Cw^{-2p-2}$, uniformly in $h,r$. For sufficiently large fixed $W_0$ it is at most half the positive lower bound, while $\tanh(t)\ge1/2$. Integration over this interval proves the claim; all omitted paired terms are nonnegative.

In the scalar lower-numerator region, replace $1\le v\le R/100$ by
$1\le v\le R/(100W_0)$. Then $w\ge W_0$, $|h|\le1$, and the same Jacobian and Student bounds apply. The integrated prediction factor is now $v I_r(w,h)\ge c_0v$. Hence the positive numerator is at least $cZ_*\log R$, since shortening the upper endpoint by a fixed factor changes the logarithm only by a constant. The bad-fit contribution is still at most
$\mathbb E|v|e^{-n^{2/3}/(12T)}=o(Z_*\log R)$ in absolute value.
Therefore the raw posterior mean at $x_*$ is at least $c\log R$, and its rare-input risk is at least $c n^{1/3}$. The event probability and zero-comparator risk are unchanged.
\end{proof}

This unit-norm version still uses input laws whose two common atoms approach one another exponentially fast. It does not satisfy a fixed uniform anti-concentration condition. The conclusion is a distribution-uniform failure, not inconsistency for a single fixed input law.

\begin{corollary}[Smooth atomless input distributions]
The input laws in Corollary~\ref{cor:raw-student-circle} can instead
have $C^\infty$ densities with respect to arc length on the unit
circle. Labels remain noiseless and bounded. The same risk lower
bound, up to a fixed constant factor, and the same zero-comparator
risk upper bound hold. No uniform bound on the input densities or
their derivatives is asserted.
\end{corollary}
\begin{proof}
Fix $n$ and its three-point construction. For every parameter vector,
$f_\theta(x)$ is continuous in $x$ on the circle. The posterior
normalizer is continuous in the training inputs and labels, by
dominated convergence with the integrable bound one. The prediction
numerator is jointly continuous in the training data and test input,
by dominated convergence with the integrable bound $|v|$.
The normalizer is strictly positive everywhere. On the compact set
of circle-valued training inputs, labels in $[-1,1]$, and a
circle-valued test input, it therefore has a positive minimum.
Consequently the raw posterior prediction is jointly continuous
and uniformly continuous on this compact set. These statements are
for each fixed $n$ and its positive temperature; no uniform modulus
in $n$ is needed.

Replace each of the three input atoms by a normalized smooth bump
supported in a sufficiently small arc around that atom, retaining
its original mixture probability. Choose the three arcs disjoint,
and retain the constant label $b,-b,0$ on the respective arcs.
Couple each perturbed observation with its original cluster center.
For sufficiently small arc radii, uniform continuity ensures that
the perturbed posterior prediction, at every test input in the rare
arc, differs from the original posterior prediction at $x_*$ by
at most one, simultaneously for every training cluster assignment.
On the original no-rare, balanced-count event, the original
prediction is at least $c n^{2/3}$. The perturbed prediction is
therefore at least $(c/2)n^{2/3}$ for all sufficiently large $n$.
The event probability is unchanged because it depends only on the
cluster assignments. The rare arc still has probability $1/n$ and
label zero, so it contributes at least $(c^2/4)n^{1/3}$ to risk.
The zero-comparator risk remains exactly $(1-1/n)b^2$.
The finite mixture of the three smooth bump densities is smooth
and atomless, proving the claim.
\end{proof}

The arc radii may depend arbitrarily strongly on $n$. This smoothing
does not give a counterexample under a fixed density bound or a
fixed quantitative anti-concentration condition.
